\documentclass[spanish]{beamer}
\usepackage[utf8]{inputenc}
\usepackage[T1]{fontenc}
\setcounter{secnumdepth}{3}
\setcounter{tocdepth}{3}
\usepackage{graphicx}
\makeatletter

%%%%%%%%%%%%%%%%%%%%%%%%%%%%%% LyX specific LaTeX commands.
%% Because html converters don't know tabularnewline
\providecommand{\tabularnewline}{\\}

%%%%%%%%%%%%%%%%%%%%%%%%%%%%%% Textclass specific LaTeX commands.
 % this default might be overridden by plain title style
 \newcommand\makebeamertitle{\frame{\maketitle}}%
 % (ERT) argument for the TOC
 \AtBeginDocument{%
   \let\origtableofcontents=\tableofcontents
   \def\tableofcontents{\@ifnextchar[{\origtableofcontents}{\gobbletableofcontents}}
   \def\gobbletableofcontents#1{\origtableofcontents}
 }

%%%%%%%%%%%%%%%%%%%%%%%%%%%%%% User specified LaTeX commands.
\usetheme{Boadilla}
\usecolortheme{whale}
\date{}

\makeatother

\usepackage{babel}
\begin{document}

\title[Métodos numéricos II]{Método de Romberg}


\author[MAC]{
Diego Hernández Landeros\\
Murillo Aceves Ana Karen\\
Fernanda Ibarra\\
Daniel Escalante\\
\begin{figure}[h]
\centering
\includegraphics[width=2cm]{mac.png}
\end{figure}
}


\institute[]{
UNAM-FES Acatlán, Matemáticas Aplicadas y Computación
}

\makebeamertitle


\begin{frame}[plain]{Contenido}


\tableofcontents{}
\end{frame}



\section{Introducción}
\begin{frame}{¿Qué es?}


La integración de Romberg es una \alert{técnica} de diseño para obtener integrales numéricas \alert{(aproximaciones)} de funciones de manera eficiente, que se basa en aplicaciones sucesivas de \alert{la regla del trapecio}. Sin embargo a través de las manipulaciones matemáticas, se alcanzan mejores resultados con menos trabajo.\\
\\$\,$\\ %SALTO DE LINEA

\end{frame}

\begin{frame}{Problema}

El problema en la práctica se presenta cuando nos vemos imposibilitados de encontrar la
función primitiva requerida, aún para integrales aparentemente sencillas como\\
{\small \[I=\int_{1}^{0}e^{x^{2}}\]}
la cual simplemente es imposible de resolver con el Teorema Fundamental del Cálculo.
\\$\,$\\
El algoritmo de Romberg forma parte de una técnica conocida como método de extrapolación de
Richardson.
\end{frame}



\section{Contenido principal}
\begin{frame}{Contenido principal}

\begin{itemize}
\item Primero de la lista
\item Segundo de la lista
\item Tercero de la lista
\end{itemize}
\end{frame}



\section{Resultados}
\begin{frame}{Marco en la sección resultados}

\begin{enumerate}
\item Primer item en la lista numerada

\begin{itemize}
\item sub item listado
\item segundo sub item listado
\end{itemize}
\item Segundo item en la lista numerada
\item Tercer item en la lista numerada
\end{enumerate}
\end{frame}



\section{Textos remarcados}
\begin{frame}{Algo}




\begin{block}{Importante} 
Texto muestra las formulas siempre van entre \textit{aquí la formula}
$\lim_{x \to \infty} \left( \frac{1}{x} \right)$
\end{block}

\begin{alertblock}{Algo importante} 
Texto muestra en rojo 
\end{alertblock}

\begin{exampleblock}{Ejemplo} 
Texto muestra en verde
$\sqrt{(x_1 - x_2)^2 + (y_1 - y_2)^2}$
\end{exampleblock}

\end{frame}



\section{Dos columnas}

\begin{frame} \frametitle{Two-column slide}
\begin{columns}
\column{0.5\textwidth} This is a text in first column. $$E=mc^2$$ \begin{itemize} \item First item \item Second item \end{itemize}
\column{0.5\textwidth} This text will be in the second column and on a second tought this is a nice looking layout in some cases. \end{columns} \end{frame}




\section{Texto}
\begin{block}{}
\begin{tabular}{|c|c|}
\hline 
\textbackslash{}textbf\{Texto muestra\} & \textbf{Texto muestra}\tabularnewline
\hline 
\textbackslash{}textit\{Texto muestra\}  & \textit{Texto muestra} \tabularnewline
\hline 
\textbackslash{}textsl\{Texto muestra\}  & \textsl{Sample Text}  \tabularnewline
\hline 
\textbackslash{}alert\{Texto muestra\}  & \alert{Texto muestra} \tabularnewline
\hline 
\textbackslash{}textrm\{Texto muestra\} & \textrm{Texto muestra}\tabularnewline
\hline 
\textbackslash{}textsf\{Texto muestra\}  & \textsf{Texto muestra} \tabularnewline
\hline 
\textbackslash{}color\{green\} Texto muestra & \color{green} Texto muestra\tabularnewline
\hline 
\textbackslash{}structure\{Texto muestra\}  & \structure{Texto muestra} \tabularnewline
\hline 
\textbackslash{}emph\{Texto muestra\} & \emph{Texto muestra}\tabularnewline
\hline 
\end{tabular}
\end{block}



\section{Referencias}
\begin{frame}{Titulo }

\begin{itemize}
\item Algo de contenido\cite{choi2000novel}\cite{litva1996digital}
\item Algo más de contenido\cite{litva1996digital}
\item Otro contenido\cite{van1988beamforming}\cite{litva1996digital,van1988beamforming}
\end{itemize}
\end{frame}



\begin{frame}{Bibliografia}


\bibliographystyle{plain}
\bibliography{bibli}

\end{frame}

\end{document}